% Options for packages loaded elsewhere
\PassOptionsToPackage{unicode}{hyperref}
\PassOptionsToPackage{hyphens}{url}
%
\documentclass[
  11pt,
  a4paper,
]{article}
\usepackage{amsmath,amssymb}
\usepackage{iftex}
\ifPDFTeX
  \usepackage[T1]{fontenc}
  \usepackage[utf8]{inputenc}
  \usepackage{textcomp} % provide euro and other symbols
\else % if luatex or xetex
  \usepackage{unicode-math} % this also loads fontspec
  \defaultfontfeatures{Scale=MatchLowercase}
  \defaultfontfeatures[\rmfamily]{Ligatures=TeX,Scale=1}
\fi
\usepackage{lmodern}
\ifPDFTeX\else
  % xetex/luatex font selection
  \setmainfont[]{Latin Modern Roman}
  \setmathfont[]{Latin Modern Math}
\fi
% Use upquote if available, for straight quotes in verbatim environments
\IfFileExists{upquote.sty}{\usepackage{upquote}}{}
\IfFileExists{microtype.sty}{% use microtype if available
  \usepackage[]{microtype}
  \UseMicrotypeSet[protrusion]{basicmath} % disable protrusion for tt fonts
}{}
\makeatletter
\@ifundefined{KOMAClassName}{% if non-KOMA class
  \IfFileExists{parskip.sty}{%
    \usepackage{parskip}
  }{% else
    \setlength{\parindent}{0pt}
    \setlength{\parskip}{6pt plus 2pt minus 1pt}}
}{% if KOMA class
  \KOMAoptions{parskip=half}}
\makeatother
\usepackage{xcolor}
\usepackage[margin=24mm]{geometry}
\usepackage{longtable,booktabs,array}
\usepackage{calc} % for calculating minipage widths
% Correct order of tables after \paragraph or \subparagraph
\usepackage{etoolbox}
\makeatletter
\patchcmd\longtable{\par}{\if@noskipsec\mbox{}\fi\par}{}{}
\makeatother
% Allow footnotes in longtable head/foot
\IfFileExists{footnotehyper.sty}{\usepackage{footnotehyper}}{\usepackage{footnote}}
\makesavenoteenv{longtable}
\setlength{\emergencystretch}{3em} % prevent overfull lines
\providecommand{\tightlist}{%
  \setlength{\itemsep}{0pt}\setlength{\parskip}{0pt}}
\setcounter{secnumdepth}{-\maxdimen} % remove section numbering
\ifLuaTeX
\usepackage[bidi=basic]{babel}
\else
\usepackage[bidi=default]{babel}
\fi
\babelprovide[main,import]{english}
\ifPDFTeX
\else
\babelfont{rm}[]{Latin Modern Roman}
\fi
% get rid of language-specific shorthands (see #6817):
\let\LanguageShortHands\languageshorthands
\def\languageshorthands#1{}
\usepackage{amsmath,amssymb}
\usepackage{needspace}
\setlength{\emergencystretch}{2em}
\raggedbottom
\clubpenalty=10000
\widowpenalty=10000
\displaywidowpenalty=10000
\AtBeginDocument{\hypersetup{pdfauthor={},pdfsubject={},pdfkeywords={}}}
\makeatletter
\renewcommand{\maketitle}{%
  \begin{center}
    {\LARGE\@title\par}
    \vspace{0.8em}
    {\normalsize\@date\par}
  \end{center}
  \vspace{0.8em}
}
\makeatother

\predisplaypenalty=10000
\postdisplaypenalty=0
\interdisplaylinepenalty=10000
\ifLuaTeX
  \usepackage{selnolig}  % disable illegal ligatures
\fi
\usepackage{bookmark}
\IfFileExists{xurl.sty}{\usepackage{xurl}}{} % add URL line breaks if available
\urlstyle{same}
\hypersetup{
  pdftitle={Real rooted Bell quotients for parallel paths and five cycle substitution},
  pdflang={en},
  hidelinks,
  pdfcreator={LaTeX via pandoc}}

\title{Real rooted Bell quotients for parallel paths and five cycle
substitution}
\author{}
\date{2 November 2026}

\begin{document}
\maketitle

\subsection{Abstract}\label{abstract}

We establish three results about the ordinary Brenti--Royle--Wagner
\(\tau\)-polynomial. First, a product inequality at the rational radius
\(R=21/20\) proves real-rootedness of the Bell edge quotient for any
number \(r\geq2\) of parallel paths. It suffices that every path length
be at least \(L\), where \(R^L\geq1+46r\). Second, the four-path family
\(\Theta(2,2,2,m)\) has a Bell edge quotient with simple strictly
negative roots for every \(m\geq2\). This follows from an exact rational
induction base and a fixed-source arrowhead invariant, whose full proof
is given here. Third, replacing one vertex of a five-cycle by any
ordinary \(\tau\)-real graph produces another ordinary \(\tau\)-real
graph whose edge quotient is real-rooted. This operation can be iterated
and permits universal ordinary edge attachment. Its proof realizes a
differential operator by a nonnegatively weighted matching graph.

The parallel-path results permit attachment along every complete
intersection because their factor graphs are triangle-free. The
substitution result permits edge attachment but makes no assertion for
larger clique intersections: its outputs can contain arbitrarily large
cliques. None of the results proves general clique-overlap closure, and
no historical novelty or proof-assistant verification is claimed.

\subsection{1 Ordinary tau reality and the Bell
normalization}\label{ordinary-tau-reality-and-the-bell-normalization}

All graphs are finite and simple. Path lengths count edges. If \(G\) has
\(n\) vertices, let
\nopagebreak
\[
P_G(x)=(-1)^n\chi_G(-x),
\]

where \(\chi_G\) is its chromatic polynomial. For the empty graph,
\(P_G=1\). Define the linear Bell transform by
\nopagebreak
\[
B((x)_j)=z^j,\qquad
(x)_j=x(x-1)\cdots(x-j+1),\qquad (x)_0=1.
\]

Equivalently,
\nopagebreak
\[
B(x^j)=T_j(z)=\sum_{i=0}^j S(j,i)z^i,
\qquad T_0=1,\qquad T_{j+1}=z(T_j+T_j'),
\]

where \(S(j,i)\) are Stirling numbers of the second kind.

In the notation of Brenti--Royle--Wagner {[}1{]}, the transforms \(S\)
and \(T\) send falling and rising factorials, respectively, to
monomials:
\nopagebreak
\[
S((x)_j)=z^j,\qquad
T\bigl(x(x+1)\cdots(x+j-1)\bigr)=z^j.
\]

Thus \(B=S\). With reflection \((\varepsilon f)(x)=f(-x)\), one has
\(B=\varepsilon T\varepsilon\), and {[}1, Proposition 5.1(d){]} gives
\nopagebreak
\[
B(P_G)(z)=(-1)^n\tau_{\mathrm{BRW}}(G,-z).\tag{1}
\]

Consequently real-rootedness of \(B(P_G)\) is precisely ordinary BRW
\(\tau\)-reality, not \(\sigma\)-reality of \(\chi_G\).
Deletion--contraction gives \(P_G=P_{G-e}+P_{G/e}\), so \(P_G\) has
nonnegative coefficients. The same is true of \(B(P_G)\). When its roots
are real, therefore, they are nonpositive.

We use the Bell product
\nopagebreak
\[
(f*g)(z)=\sum_{j\geq0}\frac{z^j}{j!}f^{(j)}(z)g^{(j)}(z).
\]

By {[}1, Propositions 4.2 and 4.4{]},
\nopagebreak
\[
B(FG)=B(F)*B(G),\tag{2}
\]

and \(f*g\) has only real nonpositive roots whenever both \(f\) and
\(g\) do.

\subsubsection{1.1 The edge attachment
interface}\label{the-edge-attachment-interface}

For a graph \(H\) with an edge, put
\nopagebreak
\[
Q_H(x)=\frac{P_H(x)}{x(x+1)}.
\]

The denominator divides \(P_H\). An ordinary attachment \(G\cup H\)
retains the common graph, makes no further identifications, and adds no
cross edges between the exclusive sides. If \(G\cap H=K_2\),
complete-cutset factorization {[}1, Proposition 1.1{]} gives
\nopagebreak
\[
P_{G\cup H}=P_GQ_H.\tag{3}
\]

Thus a real-rooted \(B(Q_H)\) with nonpositive roots supplies universal
edge attachment to every ordinary \(\tau\)-real \(G\), by (1)--(2). It
also proves that \(H\) itself is \(\tau\)-real, since
\nopagebreak
\[
P_H=x(x+1)Q_H,\qquad B(x(x+1))=z(z+2).
\]

The known empty and one-vertex overlap cases {[}1, Proposition 5.12{]}
then apply. If \(H\) is triangle-free, these cases and an edge exhaust
its complete intersections. We will use this interface for all three
results; only the triangle-free results yield every complete
intersection in this way.

For differential calculations, write
\nopagebreak
\[
D=\frac{d}{dz},\qquad \mathcal S=1+D,\qquad L=z\mathcal S.
\]

Here \(\mathcal S\) is a differential operator, distinct from the BRW
transform \(S=B\). The Bell identities are
\nopagebreak
\[
B(xF)=L B(F),\qquad
B(F(x+1))=\mathcal S B(F),\qquad
B((x+1)F)=\bigl(zD+z+1\bigr)B(F).\tag{4}
\]

\subsection{2 The quotient for any number of parallel
paths}\label{the-quotient-for-any-number-of-parallel-paths}

Let \(H=\Theta(\ell_1,\ldots,\ell_r)\) consist of \(r\geq2\) internally
vertex-disjoint paths of lengths \(\ell_i\geq2\) with the same two
distinct endpoints. Set
\nopagebreak
\[
N=\sum_{i=1}^r\ell_i,\qquad |V(H)|=N-r+2,\qquad
S_j(t)=1+t+\cdots+t^{j-1},\quad S_0(t)=0,\quad t=x+1.
\]

The geometric sum \(S_j(t)\) is unrelated to the transform \(S\). With
endpoint colors prescribed equal or distinct, the respective
internal-coloring counts for a path of length \(\ell\) are
\nopagebreak
\[
E_\ell(q)=\frac{(q-1)^\ell+(q-1)(-1)^\ell}{q},\qquad
D_\ell(q)=\frac{(q-1)^\ell-(-1)^\ell}{q}.
\]

These follow, for example, from the two-state endpoint-color recurrence;
the displayed expressions are polynomials after cancellation. Choosing
the common endpoint colors gives
\nopagebreak
\[
\chi_H(q)=q\prod_{i=1}^rE_{\ell_i}(q)
+q(q-1)\prod_{i=1}^rD_{\ell_i}(q).
\]

At \(q=1-t\),
\nopagebreak
\[
E_\ell(1-t)=(-1)^{\ell+1}tS_{\ell-1}(t),\qquad
D_\ell(1-t)=(-1)^{\ell+1}S_\ell(t).
\]

The reflection sign cancels all parity dependence, because
\(|V(H)|+\sum_i(\ell_i+1)=2N+2\) is even. After division by \(x(x+1)\),
this proves
\nopagebreak
\[
Q_H(t-1)=\prod_{i=1}^rS_{\ell_i}(t)
-t^{r-1}\prod_{i=1}^rS_{\ell_i-1}(t).\tag{5}
\]

This is a polynomial identity, including at \(t=1\).

The degree is \(d=N-r\). The first product in (5) is monic of degree
\(d\) and has next coefficient \(r\); the subtracted term is monic of
degree \(d-1\). Hence
\nopagebreak
\[
\sum_{Q_H(t_j-1)=0}t_j=-(r-1),\tag{6}
\]

with roots counted with multiplicity.

There is also a coefficientwise positivity argument. A monomial choice
in the subtracted product is indexed by \(0\leq k_i\leq\ell_i-2\), with
exponent \(r-1+\sum_i k_i\). Add one to its first \(r-1\) coordinates
and leave its last coordinate unchanged. This injects these choices,
preserving the exponent, into the monomial choices of the first product.
The image misses the constant monomial. Therefore \(Q_H(t-1)\) has
nonnegative \(t\)-coefficients and constant \(1\). In particular,
\nopagebreak
\[
Q_H(t-1)>0\quad(t\geq0),\qquad
Q_H(0)=\prod_i\ell_i-\prod_i(\ell_i-1)>0.\tag{7}
\]

Every real \(t\)-root is thus strictly negative. The same coloring
calculation and (5) also hold for a single path; this special case will
be used only as a corollary in Section 5.

\subsection{3 Long parallel paths with unrestricted path
count}\label{long-parallel-paths-with-unrestricted-path-count}

\textbf{Theorem 1.} Let \(r\geq2\), let \(\ell_i\geq2\), and set
\(R=21/20\). If
\nopagebreak
\[
\prod_{i=1}^r\frac{R^{\ell_i}+R}{R^{\ell_i}-1}<R,\tag{8}
\]

then \(B(Q_H)\) has only strictly negative real roots. Consequently
\(H\) can be ordinarily attached to any ordinary \(\tau\)-real graph
along any complete intersection while preserving ordinary
\(\tau\)-reality.

In particular, (8) holds whenever every \(\ell_i\geq L\) and
\nopagebreak
\[
R^L\geq1+46r.
\]

A sufficient minimum length is therefore
\nopagebreak
\[
L(r)=\min\{L\in\mathbb Z:L\geq2,\ 21^L\geq(1+46r)20^L\}
=\left\lceil\frac{\log(1+46r)}{\log(21/20)}\right\rceil.\tag{9}
\]

This is a sufficient bound, not an asserted optimal threshold. The cases
\(r=2\) and \(r=3\) overlap cycle and three-path results; no all-length
conclusion for arbitrary \(r\geq4\) is asserted.

\subsubsection{3.1 The product radius
bound}\label{the-product-radius-bound}

Multiply (5) by \((t-1)^r\). If \(t\) is a root with \(\rho=|t|>1\), the
factors \(t^{\ell_i}-1\) do not vanish, and rearrangement gives
\nopagebreak
\[
t=\prod_{i=1}^r
\frac{1-t^{-(\ell_i-1)}}{1-t^{-\ell_i}}.
\]

It follows that
\nopagebreak
\[
\rho\leq\prod_{i=1}^r
\frac{1+\rho^{-(\ell_i-1)}}{1-\rho^{-\ell_i}}
=\prod_{i=1}^r\frac{\rho^{\ell_i}+\rho}{\rho^{\ell_i}-1}.\tag{10}
\]

For every integer \(\ell\geq2\),
\nopagebreak
\[
f_\ell(\rho)=\frac{\rho^\ell+\rho}{\rho^\ell-1}
=1+\frac{\rho+1}{\rho^\ell-1}
\]

is strictly decreasing on \(\rho>1\): the numerator of the derivative of
its second term is
\nopagebreak
\[
\rho^\ell-1-\ell\rho^{\ell-1}(\rho+1)
=-(\ell-1)\rho^\ell-\ell\rho^{\ell-1}-1<0.
\]

If \(\rho\geq R\), (8) makes the right side of (10) less than
\(R\leq\rho\), a contradiction. Thus
\nopagebreak
\[
|t_j|<R\quad\text{for every root of }Q_H(t-1).\tag{11}
\]

\subsubsection{3.2 The precise factor repair
input}\label{the-precise-factor-repair-input}

We use two radius-disk repair lemmas from the public all-theta proof
{[}3, Section 4{]}. That reference and its complete rational
certificates are fixed at commit
\texttt{\nolinkurl{cfa1f5fd798116b44ddc95da6f90710e48fac14a}}; they are
not finite tests of the present \(r\)-path family. We state the lemmas
and explain their exact acceptance conditions.

A nonreal conjugate \(t\)-pair gives a real monic factor
\nopagebreak
\[
h(x)=x^2+px+q,\qquad p=2-2\operatorname{Re}t,\qquad q=|t-1|^2>0,
\]

and
\nopagebreak
\[
B(h)=z^2+(p+1)z+q.
\]

Call it good if \(J=(p+1)^2-4q\geq0\), and bad otherwise. A good factor
has strictly negative Bell roots because \(p+1>3-2R=9/10\). For a bad
factor, put \(\delta=q-p=|t|^2-1\). Then
\nopagebreak
\[
J=(p-1)^2-4\delta<0,\qquad
0<\delta<\frac{41}{400}<\frac19,
\qquad |p-1|<\frac23,
\qquad \operatorname{Re}t>\frac16.\tag{12}
\]

The two repair lemmas are as follows.

\begin{itemize}
\tightlist
\item
  If \(h_1\) is a bad quadratic from the radius-\(R\) disk and \(h_2\)
  is any nonreal conjugate-pair quadratic from the same disk, then
  \(B(h_1h_2)\) has four distinct strictly negative roots.
\item
  If \(h=x^2+Ax+(A+e)\), \(0\leq e\leq1/9\), \((A-1)^2\leq4e\), and
  \(1\leq\alpha\leq21/10\), then \(B(h(x)(x+\alpha))\) has three
  distinct strictly negative roots.
\end{itemize}

Here and below \(h(x)(x+\alpha)\) denotes multiplication, not
substitution into \(h\).

For the first lemma, enlarge to parameters \(a,b,u,v\in[0,1]\) and put
\nopagebreak
\[
\begin{aligned}
P&=1+\frac23b(2a-1),&Q&=P+\frac{b^2}{9},\\
p&=2+2R-4Ru,&q&=(1+R-2Ru)^2+4R^2u(1-u)v.
\end{aligned}\tag{13}
\]

These cover the bad-factor region for \(x^2+Px+Q\) and the entire radius
disk for the second quadratic. Direct expansion gives
\nopagebreak
\[
\begin{aligned}
f(z)=B\bigl((x^2+Px+Q)(x^2+px+q)\bigr)
={}&z^4+(6+P+p)z^3\\
&+\bigl(7+3(P+p)+Q+q+Pp\bigr)z^2\\
&+\bigl(1+P+p+Q+q+Pp+Pq+pQ\bigr)z+Qq.
\end{aligned}\tag{14}
\]

Its discriminant, after (13), has coordinate degrees \((10,5,6,10)\) in
\((u,v,a,b)\). All \(5,082\) tensor Bernstein coefficients are strictly
positive, with exact minimum
\nopagebreak
\[
\frac{1319869198532766581}{19683000000000000}>0.\tag{15}
\]

For precision, if a polynomial of coordinate degrees \(\mathbf n\) has
power coefficients \(c_{\mathbf j}\), the Bernstein coefficients used
here are
\nopagebreak
\[
\beta_{\mathbf i}=\sum_{\mathbf j\leq\mathbf i}
c_{\mathbf j}\prod_\ell
\frac{\binom{i_\ell}{j_\ell}}{\binom{n_\ell}{j_\ell}}.
\]

The certificate verifies positivity of every rational
\(\beta_{\mathbf i}\) and reconstructs the power coefficients by
\nopagebreak
\[
c_{\mathbf j}=
\left(\prod_\ell\binom{n_\ell}{j_\ell}\right)
\sum_{\mathbf i\leq\mathbf j}(-1)^{|\mathbf j-\mathbf i|}
\beta_{\mathbf i}\prod_\ell\binom{j_\ell}{i_\ell}.\tag{16}
\]

The tensor Bernstein basis is nonnegative and sums to one on the cube,
so (15)--(16) prove a strictly positive discriminant throughout that
cube. At \(b=0\), \(u=31/42\), \(v=300/341\), both factors equal
\(x^2+x+1\) and
\nopagebreak
\[
f(z)=z^4+8z^3+16z^2+8z+1
=\bigl(z^2+(4+\sqrt2)z+1\bigr)
 \bigl(z^2+(4-\sqrt2)z+1\bigr).
\]

This has four distinct negative roots. A continuous monic
real-polynomial family cannot change its number of real roots without a
repeated root. Connectedness of the cube and the nonvanishing
discriminant therefore give four distinct real roots everywhere in the
cube. The inequalities \(P\geq1/3\), \(Q\geq P\), \(p\geq-1/10\),
\(q\geq0\) make the nonconstant coefficients of (14) positive. Thus no
root is positive. For the actual nonreal factors, \(Qq>0\), excluding
zero as well. Boundary points with \(q=0\) need not have strictly
negative roots; they are not used as nonreal factors.

For the second lemma, direct expansion gives
\nopagebreak
\[
B(h(x)(x+\alpha))
=z^3+(3+A+\alpha)z^2
+(1+2A+\alpha+\alpha A+e)z+\alpha(A+e).\tag{17}
\]

The parameterization
\nopagebreak
\[
y=2u-1,\quad A=1+\frac23y,\quad
e=\frac{y^2+(1-y^2)v}{9},\quad
\alpha=1+\frac{11}{10}w,\qquad u,v,w\in[0,1]
\]

covers its hypotheses. The discriminant has coordinate degrees
\((6,3,4)\); all \(140\) Bernstein coefficients are positive, with exact
minimum
\nopagebreak
\[
\frac{72882116}{6834375}>0.\tag{18}
\]

The same exact conversion and inverse reconstruction are checked. At
\((u,v,w)=(1/2,0,0)\), the cubic is \((z+1)(z^2+4z+1)\). Connectedness
gives three distinct real roots throughout the cube, and its positive
coefficients make them strictly negative. The full coefficient lists and
regeneration programs for both lemmas are publicly available in {[}3{]}.
Their mathematical role is positivity on entire parameter boxes, not
extrapolation from sampled parameters.

\subsubsection{3.3 Assembly for all path
counts}\label{assembly-for-all-path-counts}

Factor \(Q_H\) over \(\mathbb R\), retaining multiplicities. By (7) and
(11), a real \(t\)-root gives a linear factor \(x+\alpha\) with
\nopagebreak
\[
1<\alpha<1+R=41/20<21/10.
\]

Its Bell transform already has a negative root. Good quadratic factors
also already have strictly negative Bell roots.

If the number of bad quadratics is even, group them in pairs and apply
the first repair lemma. If it is odd and a real linear factor exists,
combine one bad quadratic with that real factor by the second lemma and
pair the rest.

If it is odd and no real factors exist, at least one good quadratic must
exist. Otherwise every \(t\)-root would have positive real part by (12),
contradicting the negative root sum (6). Reserve one good quadratic,
combine it with one bad quadratic by the first lemma, and pair the
remaining bad quadratics.

Each factor occurrence is used once, so repeated factors present no
difficulty. All resulting groups have Bell transforms with negative real
roots. Iterating (2) proves that \(B(Q_H)\) has only real nonpositive
roots. Equation (7) gives \(B(Q_H)(0)=Q_H(0)>0\), so no zero root
occurs. The edge-attachment interface of Section 1.1 now applies. Every
cycle of \(H\) uses two paths and has length at least four, so \(H\) is
triangle-free and every complete intersection is covered.

\subsubsection{3.4 The explicit minimum length
bound}\label{the-explicit-minimum-length-bound}

If \(\ell_i\geq L\) and \(R^L\geq1+46r\), then
\nopagebreak
\[
f_{\ell_i}(R)=1+\frac{R+1}{R^{\ell_i}-1}
\leq1+\frac{41}{920r}<1+\frac{1}{22r}.
\]

Since \(\binom rk\leq r^k\), the binomial theorem gives
\nopagebreak
\[
\left(1+\frac{1}{22r}\right)^r
\leq\sum_{k=0}^r\left(\frac1{22}\right)^k
<\sum_{k=0}^\infty\left(\frac1{22}\right)^k
=\frac{22}{21}<\frac{21}{20}=R.
\]

This proves (8) and completes Theorem 1. For example, the sufficient
bounds are \(L(4)=107\), \(L(5)=112\), \(L(10)=126\), and
\(L(100)=173\). These follow from exact integer comparisons in (9); they
are not claims about sharp thresholds.

\subsection{4 Every member of the four path family Theta 2 2 2
m}\label{every-member-of-the-four-path-family-theta-2-2-2-m}

\textbf{Theorem 2.} For every integer \(m\geq2\), let
\(H_m=\Theta(2,2,2,m)\). Then \(B(Q_{H_m})\) has only simple strictly
negative roots. Ordinary attachment of \(H_m\) to any ordinary
\(\tau\)-real graph along any complete intersection preserves ordinary
\(\tau\)-reality.

\subsubsection{4.1 Exact recurrence and the fixed
source}\label{exact-recurrence-and-the-fixed-source}

Equation (5) gives, for \(m\geq1\),
\nopagebreak
\[
Q_{H_m}(t-1)=(t+1)^3S_m(t)-t^3S_{m-1}(t).
\]

Introduce the algebraic seed \(Q_0(t-1)=t^2\). This is not an assertion
about a length-zero simple graph. The sequence agrees with the displayed
graph quotient for \(m\geq1\) and satisfies
\nopagebreak
\[
Q_{m+1}(t-1)=tQ_m(t-1)+C(t),\qquad
C(t)=(t+1)^3-t^3=3t^2+3t+1.\tag{19}
\]

Indeed the first step gives \(Q_1(t-1)=(t+1)^3\), and \(S_{m+1}=tS_m+1\)
proves every later step. Thus, by (4), the polynomials \(F_m=B(Q_m)\)
obey
\nopagebreak
\[
F_0=z^2+3z+1,\qquad
F_{m+1}=zF_m'+(z+1)F_m+h,\qquad
h=3z^2+12z+7.\tag{20}
\]

In particular, \(F_m\) is monic of degree \(m+2\) and \(F_m(0)=1+7m\).

\subsubsection{4.2 A general fixed source arrowhead
invariant}\label{a-general-fixed-source-arrowhead-invariant}

The following lemma permits any fixed real polynomial source of degree
below the current polynomial. In particular it is not restricted to a
linear source.

\textbf{Lemma.} Let \(f\) be monic of degree \(n\geq2\), with distinct
strictly negative roots \(r_1,\ldots,r_n\). Let \(h\) be real,
\(\deg h<n\), and \(h(0)>0\). Define
\nopagebreak
\[
R_i=-r_i>0,\qquad a_i=\frac{h(r_i)}{f'(r_i)},\qquad
\kappa_i=\frac{a_i}{R_i},\qquad
w=\frac{h(0)}{f(0)},\qquad
V=\sum_i\frac{\kappa_i^2}{1-\kappa_i}.
\]

If
\nopagebreak
\[
|\kappa_i|<1\quad\text{for all }i,\qquad
0<w\leq\frac27,\qquad V\leq2w,\tag{21}
\]

then
\nopagebreak
\[
g=zf'+(z+1)f+h
\]

is monic of degree \(n+1\) with distinct strictly negative roots. With
the same \(h\), its new quantities satisfy
\nopagebreak
\[
|\kappa'_j|<1,\qquad
w'=\frac{w}{1+w}\leq\frac27,\qquad V'\leq2w'.\tag{22}
\]

\textbf{Proof.} Partial fractions and evaluation at zero give
\nopagebreak
\[
\frac{h(z)}{f(z)}=\sum_i\frac{a_i}{z-r_i},\qquad
\sum_i\kappa_i=w.
\]

Consequently,
\nopagebreak
\[
\frac{g(z)}{f(z)}
=z+n+1+\sum_i\frac{r_i+a_i}{z-r_i}
=z+n+1-\sum_i\frac{b_i^2}{z-r_i},\qquad
b_i=\sqrt{R_i-a_i}>0.\tag{23}
\]

The last inequality follows from \(\kappa_i<1\). Thus \(g\) is the
characteristic polynomial of the real symmetric arrowhead matrix
\nopagebreak
\[
M=\begin{pmatrix}
\operatorname{diag}(r_i)&b\\
b^{\mathsf T}&-(n+1)
\end{pmatrix}.
\]

Put \(A=-M\) and \(\Delta=\operatorname{diag}(R_i)\). The upper-left
block \(\Delta\) is positive definite; its Schur complement in \(A\) is
\nopagebreak
\[
n+1-b^{\mathsf T}\Delta^{-1}b
=n+1-\sum_i(1-\kappa_i)=1+w>0.
\]

Hence \(A\) is positive definite and every eigenvalue of \(M\) is
negative. On each interval complementary to the poles, the right side of
(23) has derivative
\nopagebreak
\[
1+\sum_i\frac{b_i^2}{(z-r_i)^2}>0.
\]

Its one-sided limits give exactly one zero on each of the \(n+1\)
complementary intervals. No pole is a zero of \(g\), because
\(g(r_i)=-b_i^2f'(r_i)\ne0\). Therefore the roots of \(g\) are distinct
and strictly negative.

To control the new residues, put \(c_i=a_i/(2b_i)\) and
\nopagebreak
\[
E=\begin{pmatrix}0&c\\c^{\mathsf T}&0\end{pmatrix}.
\]

Let \(s_j\) be a root of \(g\) and \(v_j\) a normalized real eigenvector
of \(M\) for \(s_j\). Its last coordinate is nonzero, and
\nopagebreak
\[
(v_j)_i=\frac{b_i(v_j)_{n+1}}{s_j-r_i},\qquad
(v_j)_{n+1}^2=
\frac{1}{1+\sum_i b_i^2/(s_j-r_i)^2}.
\]

Using (23) and \(g(s_j)=0\) yields
\nopagebreak
\[
v_j^{\mathsf T}Ev_j
=\frac{\sum_i a_i/(s_j-r_i)}
 {1+\sum_i b_i^2/(s_j-r_i)^2}
=\frac{h(s_j)}{g'(s_j)}.\tag{24}
\]

Now let \(K=A^{-1/2}EA^{-1/2}\). Since \(Av_j=(-s_j)v_j\), its diagonal
entries in the orthonormal eigenbasis \(v_j\) are
\nopagebreak
\[
v_j^{\mathsf T}Kv_j
=\frac{h(s_j)}{(-s_j)g'(s_j)}=\kappa'_j.\tag{25}
\]

The two nonzero eigenvalues of the rank-two matrix \(K\) are the
generalized eigenvalues of \(Ev=\lambda Av\). For \(\lambda\ne0\), the
upper-left block of \(E-\lambda A\) is \(-\lambda\Delta\). Taking its
Schur complement, and using
\nopagebreak
\[
c^{\mathsf T}\Delta^{-1}b=\frac w2,\qquad
c^{\mathsf T}\Delta^{-1}c=\frac V4,
\]

gives
\nopagebreak
\[
(1+w)\lambda^2-w\lambda-\frac V4=0.\tag{26}
\]

Indeed the Schur complement is \(-\lambda(1+w)+w+V/(4\lambda)\). Since
\(w>0\), at least one \(a_i\ne0\), so \(V>0\) and \(E\) has rank two.
Equation (26) has one positive root \(\lambda_+\) and one negative root
\(\lambda_-\). Its values at \(1\) and \(-1\) are respectively
\(1-V/4>0\) and \(1+2w-V/4>0\), because \(V\leq2w\leq4/7<4\). Hence
\nopagebreak
\[
-1<\lambda_-<0<\lambda_+<1.
\]

The remaining eigenvalues of \(K\) are zero. Every diagonal entry in any
orthonormal basis is a convex combination of its eigenvalues, so (25)
proves \(|\kappa'_j|<1\).

Finally,
\nopagebreak
\[
\phi(u)=\frac{u^2}{1-u}=\frac1{1-u}-1-u
\]

is convex for \(u<1\), since \(\phi''(u)=2/(1-u)^3>0\). Orthogonally
diagonalize \(K\) and apply scalar Jensen to each diagonal entry.
Summing, and using the row and column squared-norm identities of an
orthogonal matrix, gives
\nopagebreak
\[
V'=\sum_j\phi(\kappa'_j)
\leq\phi(\lambda_+)+\phi(\lambda_-).\tag{27}
\]

The zero eigenvalues contribute \(\phi(0)=0\). From (26),
\nopagebreak
\[
\lambda_++\lambda_-=\frac{w}{1+w},\qquad
\lambda_+\lambda_-=-\frac{V}{4(1+w)}.
\]

Thus the right side of (27) equals
\nopagebreak
\[
\Phi(V,w)=\frac{2+w}{1-V/4}-2-\frac{w}{1+w}.
\]

It is increasing in \(V<4\). Since \(V\leq2w\),
\nopagebreak
\[
V'\leq\Phi(2w,w),\qquad
\frac{2w}{1+w}-\Phi(2w,w)
=\frac{w(2-7w)}{(2-w)(1+w)}\geq0.
\]

Also \(g(0)=f(0)+h(0)\), so \(w'=w/(1+w)\leq2/7\). This gives
\(V'\leq2w'\) and proves the lemma. \(\square\)

\subsubsection{4.3 A fully specified rational induction
base}\label{a-fully-specified-rational-induction-base}

The ascending coefficient lists of \(F_0,\ldots,F_4\), reconstructed
from (20), are
\nopagebreak
\[
\begin{aligned}
F_0 &: [1,3,1],\\
F_1 &: [8,19,9,1],\\
F_2 &: [15,58,49,13,1],\\
F_3 &: [22,143,208,101,18,1],\\
F_4 &: [29,320,770,612,191,24,1].
\end{aligned}\tag{28}
\]

All terminating decimals in the next tables denote exact rational
numbers. For \(m=0,1,2,3\), the following disjoint negative intervals
have the listed endpoint signs of \(F_m\). Each row supplies exactly
\(\deg F_m\) intervals.

\begin{longtable}[]{@{}
  >{\raggedleft\arraybackslash}p{(\columnwidth - 4\tabcolsep) * \real{0.0650}}
  >{\raggedright\arraybackslash}p{(\columnwidth - 4\tabcolsep) * \real{0.6350}}
  >{\raggedright\arraybackslash}p{(\columnwidth - 4\tabcolsep) * \real{0.3000}}@{}}
\toprule\noalign{}
\begin{minipage}[b]{\linewidth}\raggedleft
\(m\)
\end{minipage} & \begin{minipage}[b]{\linewidth}\raggedright
Root intervals in increasing order
\end{minipage} & \begin{minipage}[b]{\linewidth}\raggedright
Endpoint signs in the same order
\end{minipage} \\
\midrule\noalign{}
\endhead
\bottomrule\noalign{}
\endlastfoot
\(0\) & \((-2.6181,-2.6180)\); \((-0.3820,-0.3819)\) & \((+,-)\);
\((-,+)\) \\
\(1\) & \((-6.1005,-6.1004)\); \((-2.3389,-2.3388)\);
\((-0.5607,-0.5606)\) & \((-,+)\); \((+,-)\); \((-,+)\) \\
\(2\) & \((-7.4015,-7.4014)\); \((-3.7052,-3.7051)\);
\((-1.5378,-1.5377)\); \((-0.3557,-0.3556)\) & \((+,-)\); \((-,+)\);
\((+,-)\); \((-,+)\) \\
\(3\) & \((-9.5287,-9.5286)\); \((-4.9700,-4.9699)\);
\((-2.3702,-2.3701)\); \((-0.9178,-0.9177)\); \((-0.2136,-0.2135)\) &
\((-,+)\); \((+,-)\); \((-,+)\); \((+,-)\); \((-,+)\) \\
\end{longtable}

\Needspace{160pt}
For \(f=F_4\) and \(h=3z^2+12z+7\), the exact root and residue data are:

\begin{longtable}[]{@{}
  >{\raggedright\arraybackslash}p{(\columnwidth - 4\tabcolsep) * \real{0.4300}}
  >{\raggedright\arraybackslash}p{(\columnwidth - 4\tabcolsep) * \real{0.2300}}
  >{\raggedright\arraybackslash}p{(\columnwidth - 4\tabcolsep) * \real{0.3400}}@{}}
\toprule\noalign{}
\begin{minipage}[b]{\linewidth}\raggedright
Root interval \(I_i\)
\end{minipage} & \begin{minipage}[b]{\linewidth}\raggedright
Signs of \(f\) at the endpoints
\end{minipage} & \begin{minipage}[b]{\linewidth}\raggedright
A rational enclosure for \(\kappa_i\)
\end{minipage} \\
\midrule\noalign{}
\endhead
\bottomrule\noalign{}
\endlastfoot
\((-11.62810958,-11.62810957)\) & \((+,-)\) &
\((-0.000450,-0.000449)\) \\
\((-6.60727567,-6.60727566)\) & \((-,+)\) & \((0.002892,0.002893)\) \\
\((-3.50632808,-3.50632807)\) & \((+,-)\) & \((-0.001066,-0.001065)\) \\
\((-1.58861933,-1.58861932)\) & \((-,+)\) & \((-0.019158,-0.019157)\) \\
\((-0.54542894,-0.54542893)\) & \((+,-)\) & \((-0.028264,-0.028263)\) \\
\((-0.12423842,-0.12423841)\) & \((-,+)\) & \((0.287423,0.287424)\) \\
\end{longtable}

Here is the exact certificate criterion. For a rational interval
\(I_i\), evaluate \(f\) at both endpoints using rational Horner
arithmetic and check the displayed strict opposite signs. Evaluate
\(f'\) and \(h\) on \(I_i\) by interval Horner arithmetic. Multiply the
enclosure of \(f'\) by \(-I_i\) and check that the resulting denominator
interval excludes zero before interval division. Dividing the enclosure
of \(h\) by that denominator gives a rational enclosure of
\nopagebreak
\[
\kappa_i=\frac{h(r_i)}{(-r_i)f'(r_i)}
\]

inside the last table column. All operations are additions,
multiplications and comparisons of rational numbers. The intermediate
value theorem gives a distinct negative root in each interval; degree
counting exhausts the roots and proves simplicity. Thus these tables do
not depend on numerical root classification.

The full rational enclosures, the coefficients reconstructed from (20),
and exact endpoint signs are supplied in
\href{theta222m_exact_base.json}{the base certificate};
\href{theta222m_exact_base.py}{the base checker} replays the rational
inequalities. The displayed residue enclosures are outward rational
simplifications of those full enclosures.

In particular, the first five residues have \(|\kappa_i|<1/20\) and the
last has \(|\kappa_6|<1/3\). For \(|u|\leq c<1\) one has
\(u^2/(1-u)\leq c^2/(1-c)\), so
\nopagebreak
\[
V<5\frac{(1/20)^2}{1-1/20}+\frac{(1/3)^2}{1-1/3}
=\frac{41}{228}<\frac{14}{29}=2w,\qquad
w=\frac7{29}<\frac27.\tag{29}
\]

Every hypothesis of the lemma holds at \(F_4\). Apply it repeatedly with
this fixed quadratic \(h\). The degree condition remains valid, the
roots remain simple and negative, and \(w\) decreases by
\(w\mapsto w/(1+w)\). This is an unbounded induction for every
\(m\geq4\), not an inference from a finite root sweep. The explicitly
certified smaller polynomials cover \(m=0,1,2,3\).

For graph lengths \(m\geq2\), (19)--(20) identify \(F_m=B(Q_{H_m})\).
The attachment interface applies, and \(H_m\) is triangle-free. This
completes Theorem 2.

\subsection{5 Five cycle substitution with any ordinary tau real
input}\label{five-cycle-substitution-with-any-ordinary-tau-real-input}

\textbf{Theorem 3.} Let \(F\) be any finite simple ordinary
\(\tau\)-real graph, including the empty graph. Add four new vertices
\(u,a,b,v\) inducing the path \(u-a-b-v\), and join every vertex of
\(F\) to both \(u\) and \(v\), adding no other edges. Call the resulting
graph \(H\).

Then \(B(Q_H)\) has only real nonpositive roots, \(H\) is ordinary
\(\tau\)-real, and \(H\) can be ordinarily attached along any edge to
any ordinary \(\tau\)-real graph while preserving ordinary
\(\tau\)-reality. The operation is iterable: its output again satisfies
the ordinary \(\tau\)-real input hypothesis.

Equivalently, \(H\) is obtained by replacing one vertex of a five-cycle
by \(F\), preserving its internal edges and making its two former
neighbors adjacent to every vertex of \(F\). For empty \(F\) the result
is the four-vertex path. No real-rootedness hypothesis on a quotient of
\(P_F\) is required.

\subsubsection{5.1 The positive weighted matching
operator}\label{the-positive-weighted-matching-operator}

We first prove that
\nopagebreak
\[
\mathcal A=z^2\mathcal S^4
+z\mathcal S^2(3+4D)+\mathcal S(1+3D)\tag{30}
\]

preserves the cone of real polynomials with nonnegative coefficients and
only nonpositive roots. This is a cone-preservation statement, not
preservation of all real-rooted polynomials without sign restrictions.

For \(n\geq0\), start with six core vertices \(0,1,2,3,4,5\). Its only
core edges are \(01\), \(12\), \(23\), with weights \(2\), \(1/2\),
\(1/2\), respectively; vertices \(4,5\) have no core edges. Add \(n\)
ordinary pairs \((r_i,r'_i)\). Each \(r_i\) is joined to \(2\) with
weight \(1/2\). Each \(r'_i\) is joined to \(1,2,4,5\) with weights
\(1,1/2,1,1\), respectively. There are initially no diagonal edges
\(r_ir'_i\) and no further edges.

Let \(m_j\) be the sum of products of edge weights over all \(j\)-edge
matchings. Define the reverse matching polynomial
\nopagebreak
\[
R_n(z)=\sum_jm_jz^{n+2-j}.
\]

This is a polynomial: \(\{r'_1,\ldots,r'_n,1,2\}\) is a vertex cover of
size \(n+2\), so every matching has size at most \(n+2\).

We now count all matchings for arbitrary \(n\). The possible core
matchings \(M\) are the empty matching, one of the three core edges, or
the pair \(\{01,23\}\). Their weights are \(1,2,1/2,1/2,1\). Only
support vertices
\nopagebreak
\[
U=\{1,2,4,5\}
\]

receive ordinary attachments. Set \(a_1=a_4=a_5=0\), \(a_2=1/2\), and
\(b_s=1-a_s\) for \(s\in U\). For disjoint subsets
\(A,B\subseteq U\setminus V(M)\), choose attachments using unprimed
ordinary vertices at \(A\) and primed ordinary vertices at \(B\). There
are independently \((n)_{|A|}\) and \((n)_{|B|}\) injective choices of
the two endpoint types. Using \(r_i\) and \(r'_i\) together is allowed,
since no diagonal is present. Thus the exact contribution to
\(m_{|M|+|A|+|B|}\) is
\nopagebreak
\[
\operatorname{wt}(M)(n)_{|A|}(n)_{|B|}
\prod_{s\in A}a_s\prod_{s\in B}b_s.\tag{31}
\]

For an explicit short expansion, let \(u\) count matching edges, and let
\(a,b\) mark the two attachment types before assigning their ordinary
vertices. Conditioning on the five core matchings gives
\nopagebreak
\[
\begin{aligned}
\mathcal F(u,a,b)={}&(1+bu)^3\left(1+\frac{(a+b)u}{2}\right)\\
&+2u(1+bu)^2\left(1+\frac{(a+b)u}{2}\right)\\
&+\frac u2(1+bu)^2+\frac u2(1+bu)^3
+u^2(1+bu)^2.
\end{aligned}\tag{32}
\]

Replace every \(a^ib^j\) by \((n)_i(n)_j\) and read the coefficients of
\(u\). This is precisely the finite formula (31), and expansion yields
\nopagebreak
\[
\begin{aligned}
m_0&=1,&m_1&=4n+3,\\
m_2&=6n^2+4n+1,&m_3&=n(4n^2-n+1),\\
m_4&=n(n-1)(n^2-n+1),&m_j&=0\quad(j\geq5).
\end{aligned}\tag{33}
\]

The formal falling-factorial calculation proves these identities for all
\(n\), not for a sample of degrees. Normal ordering (30) gives
coefficients of \(D^0,\ldots,D^4\) equal to
\nopagebreak
\[
z^2+3z+1,\quad
4z^2+10z+4,\quad
6z^2+11z+3,\quad
4z^2+4z,\quad z^2.
\]

Applying this operator to \(z^n\) and collecting powers produces exactly
(33). Hence
\nopagebreak
\[
R_n=\mathcal A(z^n)\qquad(n\geq0).\tag{34}
\]

For \(n=0,1\), any terms with apparently negative exponents have zero
coefficients; the vertex-cover proof already guarantees polynomiality.
The exact symbolic expansion is also replayed by
\href{certify_short_branch_core.py}{the core identity checker}, with
\href{certify_short_branch_core.json}{its rational identity record}.

\subsubsection{5.2 Diagonal edge lifting and Heilmann Lieb
normalization}\label{diagonal-edge-lifting-and-heilmann-lieb-normalization}

Let
\nopagebreak
\[
f(z)=\prod_{i=1}^n(z+\alpha_i),\qquad \alpha_i\geq0.
\]

Add a diagonal edge \(r_ir'_i\) of weight \(\alpha_i\) for each \(i\).
The same vertex cover remains valid. Classify a matching by its exact
subset \(J\) of used diagonal edges. Remove the endpoints of the
selected diagonals. Every unselected diagonal is forbidden within this
class, so the remaining allowed non-diagonal graph is exactly the
baseline graph with \(n-|J|\) ordinary pairs.

Selecting \(|J|\) diagonals reduces the reverse exponent by \(|J|\).
Thus the reverse matching polynomial of the weighted graph is
\nopagebreak
\[
\begin{aligned}
R_G(z)
&=\sum_{J\subseteq[n]}\left(\prod_{i\in J}\alpha_i\right)R_{n-|J|}(z)\\
&=\mathcal A\left[\sum_{J\subseteq[n]}
\left(\prod_{i\in J}\alpha_i\right)z^{n-|J|}\right]
=\mathcal A f(z).
\end{aligned}\tag{35}
\]

All weights are nonnegative. The weighted Heilmann--Lieb theorem {[}2{]}
applies to general graphs, not only bipartite graphs: with \(m_j\) the
nonnegative weighted matching counts,
\nopagebreak
\[
\mu_G(y)=\sum_j(-1)^jm_jy^{|V(G)|-2j}
\]

has only real roots. The auxiliary graph here need not be bipartite;
diagonal edges can create triangles. Since \(|V(G)|=2n+6\), the
normalization is exactly
\nopagebreak
\[
\mu_G(y)=(-1)^{n+2}y^2R_G(-y^2).\tag{36}
\]

If a nonzero root \(\zeta\) of \(R_G\) were not a negative real number,
a square root of \(-\zeta\) would be a nonreal zero of \(\mu_G\),
contradicting Heilmann--Lieb. Therefore all roots of
\(R_G=\mathcal A f\) are real and nonpositive; a zero root is allowed.
Repeated \(\alpha_i\), zero \(\alpha_i\), and \(n=0\) present no
difficulty. Nonnegative leading scaling and the zero polynomial complete
the stated cone-preservation result.

\subsubsection{5.3 The exact coloring identity for
substitution}\label{the-exact-coloring-identity-for-substitution}

Condition on the colors of \(u\) and \(v\). If distinct, there are
\(q(q-1)\) ordered choices; \(F\) must avoid both colors and contributes
\(\chi_F(q-2)\). The path of length three contributes \(D_3(q)\). If
equal, there are \(q\) choices, and the two factors are \(\chi_F(q-1)\)
and \(E_3(q)\). Hence
\nopagebreak
\[
\chi_H(q)=q(q-1)D_3(q)\chi_F(q-2)
+qE_3(q)\chi_F(q-1),\tag{37}
\]

where
\nopagebreak
\[
q(q-1)D_3(q)=(q-1)\bigl((q-1)^3+1\bigr),\qquad
qE_3(q)=(q-1)^3-(q-1).
\]

Since \(|V(H)|=|V(F)|+4\), unsigned reflection gives
\nopagebreak
\[
P_H(x)=x(x+1)\left[
(x^2+3x+3)P_F(x+2)-(x+2)P_F(x+1)\right].\tag{38}
\]

Using (4),
\nopagebreak
\[
B(Q_H)=\left[(L^2+3L+3)\mathcal S^2-(L+2)\mathcal S\right]B(P_F)
=\mathcal A B(P_F).\tag{39}
\]

The last equality follows by normal ordering, or by the displayed
derivative coefficients after (33). Ordinary \(\tau\)-reality of \(F\)
places \(B(P_F)\) directly in the nonnegative-coefficient,
nonpositive-root cone. Section 5.2 therefore proves that \(B(Q_H)\) has
only nonpositive real roots. No division of \(P_F\) by \(x\) or by any
clique factor occurs. In particular the empty input \(P_F=1\) is
covered.

By Section 1.1, \(H\) is ordinary \(\tau\)-real and supplies universal
ordinary edge attachment. Since \(H\) itself is an allowed input, the
substitution can be iterated. This proves Theorem 3.

\subsubsection{5.4 The short parallel path corollary and the clique
limit}\label{the-short-parallel-path-corollary-and-the-clique-limit}

Take \(F\) to be an edgeless graph on \(k\geq0\) vertices. Then
\(P_F=x^k\), \(B(P_F)=T_k\), and the output consists of \(k\) length-two
branches and one length-three branch. Its quotient is
\nopagebreak
\[
Q_k(x)=(x+2)^k(x^2+3x+3)-(x+1)^k(x+2),\qquad
B(Q_k)=\mathcal A T_k.\tag{40}
\]

For \(k=0\) this is the ordinary path of length three, not a generalized
theta graph; for \(k=1\) it is \(C_5\). For \(k\geq2\) it is
\(\Theta(2^k,3)\) with \(r=k+1\) parallel paths. This is a corollary of
Theorem 3, not a separate hypothesis or a fourth result. Its Bell
quotient has strictly negative roots, because
\nopagebreak
\[
B(Q_k)(0)=Q_k(0)=3\cdot2^k-2>0.
\]

These factor graphs are triangle-free, so their complete intersections
are exhausted by the cases in Section 1.1.

For general \(F\), however, the output can contain triangles or
arbitrarily large cliques: for example, the vertices of a clique in
\(F\), together with \(u\), still form a clique. Theorem 3 asserts
universal \textbf{edge} attachment, not attachment along cliques of size
greater than two. It is not a general graph-substitution closure
theorem.

\subsection{6 Research boundary and exact
verification}\label{research-boundary-and-exact-verification}

The three results use distinct sufficient mechanisms. Theorem 1 requires
a sufficiently large minimum path length, not merely a large maximum.
Theorem 2 proves every \(m\geq2\) for one four-path family by an
invariant. Theorem 3 allows arbitrary ordinary \(\tau\)-real
substitution inputs, but supplies only the edge-quotient attachment
interface for its potentially non-triangle-free outputs.

The radius restriction in Theorem 1 is a limitation of that proof
method. For \(\Theta(2,2,2,m)\) one has
\nopagebreak
\[
(t-1)Q_m(t-1)=t^m A(t)-C(t),\qquad
A(t)=t^3+2t^2+3t+1,\qquad C(t)=3t^2+3t+1.
\]

The cubic \(A\) has one real root \(\alpha\in(-1/2,-2/5)\), since
\(A'(t)=3t^2+4t+3>0\) on \(\mathbb R\). Its other roots are a simple
conjugate pair \(\zeta,\overline\zeta\), with \(\alpha|\zeta|^2=-1\), so
\(|\zeta|>\sqrt2>21/20\). On a sufficiently small circle about
\(\zeta\), lying in \(|t|>\sqrt2\) and excluding \(1\), the term
\(t^mA(t)\) eventually dominates \(C(t)\). Rouché's theorem supplies a
quotient root there; shrinking the circle shows convergence to
\(\zeta\). Thus quotient roots persist outside the radius-\(21/20\) disk
as \(m\) grows. This obstructs a maximum-length-only disk argument, not
Bell real-rootedness: Theorem 2 already proves that family's Bell
conclusion for all \(m\geq2\).

A broader unresolved operator family, with the differential notation of
Section 1, is
\nopagebreak
\[
\mathcal A_1=\mathcal S^2,\qquad
\mathcal A_{m+1}=L_1\mathcal A_m+D\mathcal S,\qquad
L_1=z\mathcal S+1.\tag{41}
\]

It gives the Bell quotient \(\mathcal A_mT_k\) for \(k\) length-two
branches and one length-\(m\) branch. The operator in Theorem 3 is
\(\mathcal A_3\). No preservation theorem for every \(\mathcal A_m\) is
asserted, and no finite diagnostic calculations are premises of the
three proofs above. A residue estimate uniform in arbitrary branch count
is likewise not supplied by the fixed four-path induction.

The accompanying verification programs reconstruct the exact arithmetic
obligations: the rational induction base, the formal matching identity,
the coloring/operator identities, and the two Bernstein certificates.
Section 3.2 gives the exact coefficient acceptance condition for the
latter; Section 4.3 gives the exact sign and interval criterion for the
induction base. The infinite ranges are justified by the analytic radius
argument, the invariant induction, and the formal all-degree matching
construction, not by a finite sweep. Package execution and integrity
instructions are provided separately.

AI systems assisted with derivations, exact computations, review passes,
and manuscript preparation. AI review is not independent human peer
review, and this work is ordinary computer-assisted mathematics, not
proof-assistant formalization. The full \(\tau\) clause of BRW
Conjecture 6.5 for arbitrary pairs of \(\tau\)-real graphs remains
outside the claims.

\subsubsection{6.1 Unresolved literature
overlap}\label{unresolved-literature-overlap}

A potentially relevant paper is Lili Mu, Qi Shen and Xinyao Xu, \emph{On
the Real-rootedness, invariance, and root bounds of
\(\tau\)-polynomials} {[}4{]}. Its full text, exact normalization,
hypotheses, graph families and theorem statements have not been
inspected. The accessible formula-stripped metadata does not determine
whether its results subsume, overlap with, or differ from the present
statements. That comparison remains unknown. Inaccessibility is not
evidence of non-overlap, and no historical priority or comprehensive
literature certification is claimed.

\subsection{References}\label{references}

\begin{enumerate}
\def\labelenumi{\arabic{enumi}.}
\item
  Francesco Brenti, Gordon F. Royle, and David G. Wagner. \emph{Location
  of zeros of chromatic and related polynomials of graphs}. Canadian
  Journal of Mathematics \textbf{46} (1), 55--80 (1994).
  \href{https://doi.org/10.4153/CJM-1994-002-3}{doi:10.4153/CJM-1994-002-3}.
  Proposition 1.1 gives complete-cutset factorization; Propositions 4.2
  and 4.4 give Bell-product identity and preservation; Proposition
  5.1(d) gives reflection; Proposition 5.12 covers empty and one-vertex
  intersections.
\item
  \begin{samepage}
  Ole J. Heilmann and Elliott H. Lieb. \emph{Theory of monomer-dimer
  systems}. Communications in Mathematical Physics \textbf{25}, 190--232
  (1972).
  \href{https://doi.org/10.1007/BF01877590}{doi:10.1007/BF01877590}. The
  nonnegative weighted matching theorem applies to general graphs; (36)
  records the normalization used here.

  \end{samepage}
\item
  \emph{Real rooted edge quotients and clique attachment for theta
  graphs}. Public proof and exact certificates, fixed GitHub commit
  \texttt{\nolinkurl{cfa1f5fd798116b44ddc95da6f90710e48fac14a}}, 2
  November 2026.
  \href{https://github.com/Aurantium-47/brw-tau-clique-attachments/blob/cfa1f5fd798116b44ddc95da6f90710e48fac14a/extensions/2026-11-02-all-theta/MANUSCRIPT.md\#4-exact-factor-repair-in-the-radius-disk}{Manuscript
  Section 4};
  \href{https://github.com/Aurantium-47/brw-tau-clique-attachments/tree/cfa1f5fd798116b44ddc95da6f90710e48fac14a/extensions/2026-11-02-all-theta}{immutable
  certificate package}. In particular, the
  \href{https://github.com/Aurantium-47/brw-tau-clique-attachments/blob/cfa1f5fd798116b44ddc95da6f90710e48fac14a/extensions/2026-11-02-all-theta/effective_mixed_certificate.json}{mixed
  quartic certificate} and
  \href{https://github.com/Aurantium-47/brw-tau-clique-attachments/blob/cfa1f5fd798116b44ddc95da6f90710e48fac14a/extensions/2026-11-02-all-theta/effective_cubic_certificate.json}{cubic
  certificate} supply the exact coefficient lists used in Section 3.2.
\item
  Lili Mu, Qi Shen, and Xinyao Xu. \emph{On the Real-rootedness,
  invariance, and root bounds of \(\tau\)-polynomials}. Discrete
  Mathematics, Algorithms and Applications, article 2550185.
  \href{https://doi.org/10.1142/S179383092550185X}{doi:10.1142/S179383092550185X}.
  Its precise mathematical overlap is unverified, as explained in
  Section 6.1.
\end{enumerate}

\end{document}
